\chapter*{Introduction générale}
\addcontentsline{toc}{chapter}{Introduction générale}

Dans un contexte numérique marqué par l'évolution rapide des cybermenaces, la sécurité des applications constitue un enjeu stratégique pour les organisations qui développent ou exploitent des logiciels critiques. Les systèmes liés au paiement électronique sont particulièrement concernés, car ils manipulent des données sensibles et doivent répondre à des exigences strictes de confidentialité, d'intégrité, de disponibilité et de traçabilité.

Le \textit{PCI Secure Software Framework} (PCI SSF) fournit un cadre de référence destiné à renforcer la sécurité des logiciels de paiement et des processus de développement associés. Parallèlement, les approches DevSecOps visent à intégrer la sécurité dès les premières phases du cycle de vie logiciel, en s'appuyant sur l'automatisation, la collaboration et l'amélioration continue.

Ce rapport présente le travail réalisé dans le cadre d'un stage de fin d'études au sein de Hightech Payment Systems (HPS). Il s'intéresse à la conception et à la mise en oeuvre d'une démarche de conformité PCI SSF intégrée à un cycle DevSecOps. Le document est organisé comme suit :

\begin{itemize}
    \item le premier chapitre présente l'organisme d'accueil, son organisation et l'équipe d'affectation ;
    \item le deuxième chapitre expose le contexte du projet, la problématique, les objectifs et la méthodologie ;
    \item le troisième chapitre introduit les concepts nécessaires autour du PCI SSF, du DevSecOps et de la sécurité logicielle ;
    \item le quatrième chapitre décrit la démarche proposée pour intégrer les exigences PCI SSF au cycle DevSecOps ;
    \item le cinquième chapitre présente la réalisation, l'évaluation et les résultats obtenus.
\end{itemize}

\clearpage
